%------------------------------------------------------------------------------%
\section{Integração por Substituição Trigonométrica}
\label{sec:Integracao_substituicao_Trigonometrica}

Quando tentamos resolver a integral
\[
  \int x \sqrt{a^2-x^2}dx,
\]
podemos usar uma substituição simples $u=a^2-x^2$.

No geral fazemos a seguinte mudança
\[
  \int f(g(x))g'(x) dx= \int f(u)du
\]

Para a integral
\[
  \int \sqrt{a^2-x^2}dx
\]
não conseguimos fazer a substituição simples

Nesse caso fazemos uma substituição ``inversa''
$x=a\sin( \theta )$, assim:
\begin{align*}
  \int \sqrt{a^2-x^2}\,dx
  & = \int \sqrt{a^2-a^2 \sin^2(\theta)\,} a\cos(\theta) \,d\theta \\
  & = \int \sqrt{a^2\big(1-\sin^2(\theta)\big)\,}   \cos(\theta) \,d\theta \\
  & = \int a^2 \cos^2 (\theta) d\theta
\end{align*}

Nesse caso geral temos
\[
  \int f(x)dx= \int f\big(g(\theta)\big) g'(\theta) \,d\theta
\]

Podemos fazer a substituição inversa desde que $g$ seja injetiva.

No nosso caso  podemos fazer a substituição inversa $x=a\sin(\theta)$,
desde que essa função seja injetora. Isso pode ser conseguido se
restringirmos $\theta$ ao intervalo
\[
  \left[-\frac{\pi}{2}, \frac{\pi}{2}\right]
\]

Nessa seção estudaremos a técnica de substituição trigonométrica e ela serve
para integrar funções que tem uma ``raiz'' e não conseguimos usar nenhuma das
técnicas estudadas até aqui. A Tabela~\ref{tab:substituicoes_trigonometricas}
apresenta os casos onde usaremos a substituição trigonométrica.

%------------------------------------------------------------------------------%
\begin{table}
  \centering
  \renewcommand*{\arraystretch}{2.2}
  \caption{Substituições Trigonométricas}
  \label{tab:substituicoes_trigonometricas}
  \begin{ceqn}
  \[
    \begin{array}{>{\;}c<{\;}>{\;}c<{\;}>{\;}l<{\;}>{\;}c<{\;}}
      \hline
      \text{Expressão}    &
      \text{Substituição} & &
      \text{Identidade}
      \\ \hline
      \sqrt{a^2-x^2}
      & x=a\sin(\theta)
      & -\sfrac{\pi}{2} \leq \theta \leq \sfrac{\pi}{2}
      & 1-\sin^2(\theta)= \cos^2(\theta)
      \\
      \sqrt{a^2+x^2}
      & x=a\tan(\theta)
      & -\sfrac{\pi}{2} < \theta < \sfrac{\pi}{2}
      & 1+\tan^2(\theta) = \sec^2(\theta)
      \\
      \sqrt{x^2-a^2}
      & x=a\sec(\theta)
      &
      \hpm
      \begin{array}{l}
          0 \leq \theta < \sfrac{\pi}{2} \\
        \pi \leq \theta < \sfrac{3\pi}{2}
      \end{array}
      & \sec^2(\theta)-1= \tan^2(\theta)
      \\ \hline
    \end{array}
  \]
  \end{ceqn}
\end{table}
%------------------------------------------------------------------------------%

%------------------------------------------------------------------------------%
\begin{example}
  Calcule \(\int \frac{\sqrt{9-x^2}}{x^2}dx\)
  \solution
  Note que aqui temos uma integral que envolve a raiz
  \[
    \sqrt{9-x^2} = \sqrt{3^2-x^2}
  \]
  então vamos fazer a substituição
  $x=a \sin(\theta)$, onde $a=3$.

  Seja $x= 3\sin(\theta)$, com
  \(-\frac{\pi}{2} \leq x \leq \frac{\pi}{2}\).
  Então $dx=3 \cos(\theta) d\theta$.
  Assim,
  \begin{align*}
    I
    & = \int \frac{\sqrt{9-x^2}}{x^2}dx \\
    & = \int \frac{\sqrt{9-9\sin^2(\theta)}}{9 \sin^2(\theta)} 3 \cos(\theta) d\theta \\
    & = \frac{1}{3}\int \frac{\sqrt{9(1-\sin^2(\theta))}}{ \sin^2(\theta)} \cos(\theta) d\theta
  \end{align*}
  Usando agora que
  \(1-\sin^2(x)= \cos^2(x)\)
  \begin{align*}
    I
    & = \frac{1}{3}\int \frac{\sqrt{9\cos^2(x)}}{ \sin^2(\theta)} \cos(\theta) d\theta \\
    & = \frac{1}{3}\int \frac{3\cos(\theta)}{\sin^2(\theta)}\cos(\theta) d\theta  \\
    & = \int \frac{\cos^2(\theta)}{\sin^2(\theta)} d\theta \\
    & = \int \cot^2(\theta) d\theta
  \end{align*}
  Usando
  \(\csc^2(\theta)-1 = \cot^2(\theta)\)
  temos
  \begin{align*}
  I
    & =  \int (\csc^2(\theta)-1) d \theta  \\
    & = -\cot(\theta)- \theta +C
  \end{align*}
  Agora precisamos voltar para a variável original $x$, ou seja, é necessário
  expressar $\cot(\theta)$ e $\theta$ em função de $x$. Ora, temos que
  \[
    x = 3\sin(\theta)
    \Rightarrow
    \sin(\theta) = \frac{x}{3}.
  \]
  Se interpretarmos
  $\theta$ como um ângulo de um triângulo retângulo então podemos escolher $x$
  como o cateto oposto e $3$ como hipotenusa, uma vez que
  \(\sin(\theta)=\frac{x}{3}\).
  O Teorema de Pitágoras nos dá que o cateto adjacente é
  $\sqrt{9-x^2}$. Assim,
  \[
    \cot(\theta)= \frac{\sqrt{9-x^2}}{x}
  \]

  Embora no diagrama $\theta >0$, essa expressão para $\cot(\theta)$ também é
  válida para $\theta <0$. Uma vez que \(\sin(\theta)= \frac{x}{3}\) e
  $y=\sin(\theta)$ é injetiva no intervalo
  \(-\frac{\pi}{2} \leq \theta \leq\frac{\pi}{2}\)
  então \(\theta = \arcsin\left(\frac{x}{3}\right)\), logo
  \[
    \int \frac{\sqrt{9-x^2}}{x^2}dx
    = -\cot(\theta)-\theta+C
    = \frac{\sqrt{9-x^2}}{x}-\arcsin\left(\frac{x}{3}\right)+C
  \]

  \IncludeGraphicCentred{exemplo_triangulo_1}[0.4][-1]
\end{example}
%------------------------------------------------------------------------------%

Durante a resolução do exercício usamos que
\(\sqrt{\cos^2(\theta)}=\cos(\theta)\),
fizemos isso porque
\(\theta \in \left[-\frac{\pi}{2}, \frac{\pi}{2} \right]\)
e $\cos(x)$ é não-negativo nesse intervalo, assim,
$\abs{\cos(x)}= \cos(x)$, nesse caso.

%------------------------------------------------------------------------------%
\begin{example}
  Encontre \(\int \frac{1}{x^2\sqrt{x^2+4}}dx\)
  \solution
  Fazendo a substituição trigonométrica
  $x=2\tan(\theta)$, com $-\frac{\pi}{2}< \theta < \frac{\pi}{2}$,
  temos que $dx=2\sec^2(\theta)d \theta$.
  Assim,
  \begin{align*}
    I
    & = \int \frac{1}{x^2\sqrt{x^2+4}}dx \\
    & = \int \frac{1}{\tan^2(\theta) \sqrt{4\tan^2(\theta)+4}}2\sec^2(\theta) d \theta \\
    & = \frac{1}{2}\int \frac{\sec^2(\theta) }{\tan^2(\theta) \sqrt{4(1+ \tan^2(\theta))}} d\theta
  \end{align*}
  Usando \(1+\tan^2(\theta)= \sec^2(\theta)\)
  \begin{align*}
    I
    & = \frac{1}{2} \int \frac{\sec^2(\theta)}{\tan^2(\theta) \sqrt{4\sec^2(\theta)}} d \theta \\
    & = \frac{1}{2} \int \frac{\sec^2(\theta)}{\tan^2(\theta) 2\sec(\theta)}  d \theta \\
    & = \frac{1}{4} \int \frac{\sec(\theta)}{\tan^2(\theta)} d \theta \\
    & = \frac{1}{4} \int \frac{\;\;\dfrac{1}{\cos(\theta)}\;\;}{\dfrac{\sin^2(\theta)}{\cos^2(\theta)}} d\theta \\
    & = \frac{1}{4} \int \frac{\cos(\theta)}{\sin^2(\theta)} d \theta
  \end{align*}
  Usando a substituição $u=\sin(\theta)$ temos que
  $du=\cos(\theta) d\theta$, logo
  \begin{align*}
    \frac{1}{4} \int \frac{\cos(\theta)}{\sin^2(\theta)} d \theta
    & = \frac{1}{4} \int \frac{du}{u^2} \\
    & = \frac{1}{4} - \frac{1}{u}+C \\
    & = \frac{1}{4\sin(\theta)} + C \\
    & = \frac{-\csc(\theta)}{4}+ C
  \end{align*}
  Usando a informação que $x= 2 \tan(\theta) \Rightarrow \tan(\theta)= \frac{x}{2}$
  e olhando o diagrama
  \IncludeGraphicCentred{exemplo_triangulo_2}[0.4][-1][-1]
  temos que $\csc(\theta)=\frac{x^2+4}{x}$.
  Assim,
  \[
    \int \frac{dx}{x^2\sqrt{x^2+4}} =- \frac{\sqrt{x^2+4}}{4x}+C.
  \]
\end{example}
%------------------------------------------------------------------------------%

%------------------------------------------------------------------------------%
\begin{example}
  Encontre \(\int \frac{dx}{\sqrt{x^2-a^2}}\), onde $a>0$.
  \solution
  Aqui usamos a substituição trigonométrica
  $x=a\sec(\theta ) \Rightarrow dx= a \sec(\theta) \tan(\theta) d \theta$,
  onde $0 < \theta < \frac{\pi}{2}$ ou $\pi < \theta  < \frac{3\pi}{2}$.
  Assim,
  \begin{align*}
    I
    & = \int \frac{dx}{\sqrt{x^2-a^2}} \\
    & = \int  \frac{ a \sec(\theta) \tan(\theta) d \theta}{\sqrt{a^2 \sec^2(\theta)}-a^2} \\
    & = a\int \frac{ \sec(\theta) \tan(\theta) d \theta}{\sqrt{a^2(\sec^2(\theta)-1)}}
  \end{align*}
  Usando \(\sec^2(\theta)-1= \tan^2(\theta)\)
  \begin{align*}
    I
    & = a \int \frac{\sec(\theta) \tan(\theta) d \theta}{a \sqrt{\tan^2(\theta)}} \\
    & = \int \frac{\sec(\theta)\tan(\theta) d \theta}{\tan(\theta)}= \int \sec(\theta) d \theta \\
    & = \ln \abs{\sec(\theta)+ \tan(\theta)} + C
  \end{align*}
  Já temos que $\sec(\theta)= \frac{x}{a}$ e usando o triângulo da figura
  \IncludeGraphicCentred{exemplo_triangulo_3}[0.4][-1][-1]
  temos que $\tan(\theta)=\frac{\sqrt{x^2-a^2}}{a}$.
  Logo,
  \[
    \int \frac{dx}{\sqrt{x^2-a^2}}
    = \ln \abs{\frac{x}{a}+ \frac{\sqrt{x^2-a^2}}{a}}+ C
  \]
\end{example}
%------------------------------------------------------------------------------%

Durante a resolução do exercício usamos o fato que
$\int \sec(x) dx = \ln \abs{\sec(x)+\tan(x)}+ C$,
para demonstrar esse fato multiplicamos o numerador e o denominador
por $\sec(x)+\tan(x)$, assim
\[
  \int \sec(x) dx
  = \int \sec(x) \frac{\sec(x)+\tan(x)}{\sec(x)+\tan(x)}dx
  = \int \frac{\sec^2(x)+ \sec(x)\tan(x)}{\sec(x)+\tan(x)} dx
\]
Usando a substituição $u=\sec(x)+\tan(x)$
temos que $du= (\sec(x)\tan(x)+ \sec^2(x))dx$.
Então,
\[
  \int \sec(x) dx
  = \int \frac{1}{u} du
  = \ln\abs{u}+C
  = \ln \abs{\sec(x)+\tan(x)}+ C
\]

O próximo exemplo nos mostra que mesmo quando as substituições
trigonométricas são possíveis, elas nem sempre nos dão a solução mais fácil.
Você deve primeiro procurar um método mais simples.

%------------------------------------------------------------------------------%
\begin{example}
  Encontre \(\int \frac{x}{\sqrt{x^2+4}} dx\)
  \solution
  Aqui poderíamos pensar em usar a substituição trigonométrica
  $x=2 \tan(\theta)$ (e até pode), no entanto, fazer a substituição
  simples $u=x^2+4$ é mais fácil, uma vez que $du=2x dx$ e
  \[
    \int \frac{x}{\sqrt{x^2+4}} dx
    = \frac{1}{2} \int \frac{1}{\sqrt{u}} du
    = \sqrt{u}+ C
    = \sqrt{x^2+4}+C
  \]
\end{example}
%------------------------------------------------------------------------------%

Muitas vezes precisamos usar mais de uma técnica ou a mesma técnica
de integração para resolver uma integral. O próximo exemplo ilustra
uma situação dessa.

%------------------------------------------------------------------------------%
\begin{example}
  Calcule \(\int x \arcsin(x) dx\)
  \solution
  Vamos começar resolvendo a integral usando integração por partes.
  Chame $u=\arcsin(x)$ e $dv=xdx$, de modo que
  $du=\frac{1}{\sqrt{1-x^2}}dx$ e $v=\frac{x^2}{2}$.
  Assim, usando integração por partes,
  \begin{equation}
    \label{eq:exemplo_substtri}
    \int x \arcsin(x) dx
    = \frac{x^2 \arcsin(x)}{2}-\frac{1}{2}\int \frac{x^2}{\sqrt{1-x^2}}dx
  \end{equation}
  Para resolvermos a integral $\int \frac{x^2}{\sqrt{1-x^2}}dx$,
  usaremos a substituição trigonométrica $x=\sin(\theta)$,
  assim $dx=\cos(\theta)d\theta$ e
  $-\frac{\pi}{2} \leq \theta \leq \frac{\pi}{2}$.
  Portanto,
  \begin{align*}
    \int \frac{x^2}{\sqrt{1-x^2}}dx
    & = \int \frac{\sin^2(\theta)\cos(\theta)}{\sqrt{1-\sin^2(\theta)}} \\
    & = \int \frac{\sin^2(\theta)\cos(\theta)}{\sqrt{\cos^2(\theta)}} \\
    & = \int \frac{\sin^2(\theta)\cos(\theta)}{\cos(\theta)} d \theta \\
    & = \int \sin^2(\theta) d \theta
  \end{align*}

  Sabendo que $\sin^2(\theta)= \frac{1}{2}(1-\cos(2 \theta))$,
  chegamos que
  \begin{align*}
    \int \sin^2(\theta) d \theta
    & = \frac{1}{2}\int(1-\cos(2 \theta)) \\
    & = \frac{1}{2}(\theta -\frac{\sin(2\theta)}{2}) \\
    & = \frac{1}{2}(\theta -\frac{2\sin(\theta)\cos(\theta)}{2})  \\
    & = \frac{1}{2}(\theta -\sin(\theta)\cos(\theta)) +C.
  \end{align*}
  Já sabemos que $x=\sin(\theta) \Rightarrow \theta=\arcsin(x)$.
  Do triângulo
  \IncludeGraphicCentred{exemplo_triangulo_4}[0.4][-1][-1]
  tiramos que $\cos(\theta)= \sqrt{1-x^2}$.
  Assim,
  \begin{align*}
    \int \frac{x^2}{\sqrt{1-x^2}}dx
    & =\frac{1}{2}(\theta -\sin(\theta)\cos(\theta)) \\
    & =\frac{1}{2}(\arcsin(x)-x\sqrt{1-x^2})+C
  \end{align*}
  Voltando para \eqref{eq:exemplo_substtri}, obtemos
  \[
    \int x \arcsin(x) dx
    = \frac{x^2\arcsin(x)}{2}-\frac{1}{4}(\arcsin(x)-x\sqrt{1-x^2})+C.
  \]
\end{example}
%------------------------------------------------------------------------------%

%------------------------------------------------------------------------------%
\begin{example}
  Calcule \(\int \frac{dt}{\sqrt{t^2-6t+13}}\)
  \solution
  Aqui não temos nenhuma expressão como aquelas da tabela de
  substituições trigonométricas. No entanto, temos uma raiz no denominador.
  Vamos manipular a expressão dentro da raiz para chegarmos em uma das
  expressões da tabela.

  Note que, completando quadrado temos
  \[
    t^2-6t+13
    =(t-3)^2-9+13
    =(t-3)^2+4.
  \]
  Fazendo a substituição trigonométrica $u=t-3 \Rightarrow du=dt$, obtemos
  \[
    \int   \frac{dt}{\sqrt{t^2-6t+13}}
    = \int \frac{dt}{\sqrt{(t-3)^2+4}}
    = \int \frac{du}{\sqrt{u^2+4}}
    = \int \frac{du}{\sqrt{u^2+2^2}}
  \]
  Agora, temos uma integral que envolve uma daquelas expressões
  que aparece na tabela de substituições trigonométricas.
  Portanto, fazendo a substituição trigonométrica
  $u=2\tan(\theta) \Rightarrow du=2\sec^2(\theta) d\theta$. Assim,
  \begin{align*}
    I
    \int \frac{du}{u^2+2^2}
    & = \int \frac{2\sec^2(\theta) }{\sqrt{4\tan^2(\theta)+4}}d\theta \\
    & = \int \frac{2\sec^2(\theta) }{2\sqrt{(\tan^2(\theta)+1)}}d\theta
  \end{align*}
  Usando \(\tan^2(\theta)+1= \sec^2(\theta)\)
  \begin{align*}
    I
    & = \int \frac{2\sec^2(\theta)}{2\sec(\theta)}d\theta \\
    & = \int \sec(\theta)d\theta \\
    & = \ln \abs{\sec(\theta)+ \tan(\theta)} + K
  \end{align*}
  Já temos que
  \[
    u=2\tan(\theta) \Rightarrow \tan(\theta)=\frac{u}{2}=\frac{t-3}{2},
  \]
  então pela figura
  \IncludeGraphicCentred{exemplo_triangulo_5}[0.4][-1][-1]
  temos que \(\sec(\theta)=\frac{2}{\sqrt{t^2-6t+13}}\). Assim,
  \[
    \int \frac{dt}{\sqrt{t^2-6t+13}}
    = \ln \left|\sec(\theta)
    =\frac{2}{\sqrt{t^2-6t+13}}+ \frac{t-3}{2} \right| +K.
  \]
\end{example}
%------------------------------------------------------------------------------%

%------------------------------------------------------------------------------%
